The attention mechanism has revolutionized artificial intelligence, particularly through the Transformer architecture in natural language processing. However, the extent to which this paradigm applies to the broader digital ecosystem remains underexplored. In this paper, we investigate the hypothesis that attention serves as a fundamental principle across both artificial intelligence and the Internet field, where digital products compete for limited human cognitive resources. We conduct a series of experiments comparing attention-based models against traditional non-attention baselines in critical Internet tasks: sequential recommendation and Click-Through Rate (CTR) prediction. Our results show that attention-based models consistently match or outperform baselines, with the AttentionMLP achieving a validation accuracy of 0.7740 on the Criteo dataset compared to 0.7512 for a standard MLP. Furthermore, we provide empirical evidence linking the mathematical focus of attention, measured by entropy, to prediction correctness; correct predictions exhibit significantly lower mean attention entropy (1.4107) than incorrect ones (1.4285). These findings suggest that attention functions as a cross-domain bottleneck, providing a unified framework for understanding information filtering in both neural networks and digital economies.
